% 原书第二篇2.5节170--179题；用户PDF页91--92，印刷页82--83。
\thermocapillaryproblem{170}{82}
两块竖直平板在润湿液体中部分浸没，
板间形成极小楔角 $\delta\varphi$，
楔棱保持竖直。
液体密度 $\rho$、表面张力 $\alpha$、
接触角 $\theta$。
求距楔棱水平距离为 $x$ 处液体上升高度 $h(x)$。
\exsolution
当楔角很小时，两板在距楔棱 $x$ 处的局部间距
$s(x)\simeq x\delta\varphi$（角度以弧度计）。
在薄狭缝中的弯月面，
沿板间截面的主曲率半径近似为
$R_c=s/(2\cos\theta)$，
另一方向曲率因楔角小而忽略。
Laplace 压差为
\[
\Delta p\simeq\frac{\alpha}{R_c}
=\frac{2\alpha\cos\theta}{x\delta\varphi}.
\]
这与静水柱压头 $\rho gh$ 平衡，故
\[
\boxed{h(x)=\frac{2\alpha\cos\theta}
{\rho g\,x\,\delta\varphi}}.
\]
\excheck 公式只适用于距离楔棱并非无限小、
板间距远大于分子尺度且弯月面可视为局部准静态的区域；
$x\to0$ 的形式发散不能当作真实无穷高度。

\thermocapillaryproblem{171}{82}
竖直向下的水射流从圆孔喷出。
在某水平截面处水柱直径为
$d=2.0\,\mathrm{mm}$；
向下相隔
$\ell=20\,\mathrm{mm}$ 的另一截面处，
直径变为前者的 $1/n$，
$n=1.5$。
求每秒流出的水体积 $Q$。
\exsolution
稳态不可压缩流满足
$A_1v_1=A_2v_2=Q$，
其中
$A_1=\pi d^2/4$、
$A_2=A_1/n^2$，故
$v_2=n^2v_1$。
沿自由射流的伯努利方程（两截面均大气压）：
\[
\frac12v_2^2=\frac12v_1^2+g\ell.
\]
代入速度比：
\[
(n^4-1)v_1^2=2g\ell
\quad\Longrightarrow\quad
v_1=\sqrt{\frac{2g\ell}{n^4-1}}.
\]
流量：
\[
\boxed{Q=\frac{\pi d^2}{4}
\sqrt{\frac{2g\ell}{n^4-1}}}.
\]
代入数值，
$\boxed{Q\simeq9.8\times10^{-7}\,\mathrm{m^3/s}
\approx1.0\,\mathrm{cm^3/s}}$。
\excheck $v_2/v_1=n^2$，
不是 $n$，因为截面积正比于直径的平方；
据此必须在能量式中出现 $n^4-1$。

\thermocapillaryproblem{172}{82}
一个水滴在空气中以恒定终端速度下落。
沿水滴竖直方向两极点相距
$h=2.3\,\mathrm{mm}$。
求上、下极点液面曲率半径的差，
假设表面张力为 $\alpha$。
\exsolution
水滴以恒速运动时，
在随水滴平动的理想近似中，
液滴内部水静压沿向下方向增加
$\rho gh$，即
$p_{\rm bottom}-p_{\rm top}=\rho gh$。
极点处外界空气压强近似相等，
由 Laplace 方程
$p_{\rm in}-p_{\rm out}=2\alpha/R_c$，
\[
2\alpha\left(\frac1{R_{\rm bottom}}
-\frac1{R_{\rm top}}\right)=\rho gh.
\]
由于水滴近似球形，
两端曲率半径均接近 $h/2$，
可在微小变形时作一阶线性化：
\[
\frac1{R_b}-\frac1{R_t}
\simeq\frac{R_t-R_b}{(h/2)^2}.
\]
故曲率半径差的模为
\[
\boxed{|R_t-R_b|\simeq\frac{\rho gh^3}{8\alpha}}.
\]
取水 $\rho\simeq1000\,\mathrm{kg/m^3}$、
$\alpha\simeq0.073\,\mathrm{N/m}$，
得到
\[
\boxed{|R_t-R_b|\simeq0.20\,\mathrm{mm}}.
\]
\excheck 这一估计要求变形较小；
恒速仅意味着水滴中心无加速度，
不等于忽略水滴内部的重力静压梯度。

\thermocapillaryproblem{173}{82}
一颗水银液滴被两块水平玻璃板夹成
近似圆片，初半径 $R$、厚度 $h\ll R$。
若在上板放置砝码使板间距减为原来的 $1/n$，
求所需砝码质量 $m$。
水银接触角为 $\theta$。
以 $R=2.0\,\mathrm{cm}$、
$h=0.38\,\mathrm{mm}$、
$n=2.0$、
$\theta=135^\circ$ 计算。
\exsolution
水银不润湿玻璃，
$\cos\theta<0$。
在薄片近似下，弯月面压力修正主导项
$|\Delta p|\simeq2\alpha|\cos\theta|/h$，
压力在圆片面积 $\pi R^2$ 上施加分离板的作用。
因为水银不可压缩，初体积
$\pi R^2h$ 为常数，
板间距变为 $h/n$ 后新半径
$R'=R\sqrt n$。
给定初末状态，外加砝码须提供的
\emph{附加压力负荷}为
\[
mg\simeq 2\pi\alpha|\cos\theta|
\left(\frac{R'^2}{h/n}-\frac{R^2}{h}\right)
=\frac{2\pi\alpha R^2|\cos\theta|}{h}(n^2-1).
\]
所以
\[
\boxed{m\simeq
\frac{2\pi\alpha R^2|\cos\theta|}{gh}(n^2-1)}.
\]
取水银 $\alpha\simeq0.48\,\mathrm{N/m}$，
得
\[
\boxed{m\simeq0.69\,\mathrm{kg}\approx0.7\,\mathrm{kg}}.
\]
\excheck 计算的是从初平衡状态额外加载的砝码质量；
必须同时考虑厚度变化引起的水银圆片半径变化。

\thermocapillaryproblem{174}{82}
两块平行玻璃板间距
$h=0.10\,\mathrm{mm}$。
在板间放入质量 $m=70\,\mathrm{mg}$
的水滴，并认为完全润湿。
求玻璃板之间的相互吸引力。
\exsolution
水在两板间形成很薄的液桥，
其主曲率造成的液内负压近似为
$|\Delta p|\simeq2\alpha/h$。
水滴体积为 $V=m/\rho$，
夹在板间形成面积
$A\simeq V/h=m/(\rho h)$。
因此负压对板的吸引力为
\[
F\simeq|\Delta p|A
=\frac{2\alpha}{h}\frac{m}{\rho h}
=\boxed{\frac{2\alpha m}{\rho h^2}}.
\]
以水 $\alpha\simeq0.073\,\mathrm{N/m}$、
$\rho=1000\,\mathrm{kg/m^3}$，
\[
\boxed{F\simeq1.0\,\mathrm N}.
\]
\excheck 需有液桥厚度远小于其横向尺寸，
使边缘表面张力的其他小修正可忽略。

\thermocapillaryproblem{175}{82}
两块半径
$R=5.0\,\mathrm{cm}$ 的玻璃圆盘被水润湿
后合拢，水层厚度为
$h=1.9\,\mathrm{\mu m}$，
假定完全润湿。
求沿圆盘法线把两盘拉开的力。
\exsolution
水层很薄，边界弯月面的主曲率
产生近似负压
$\Delta p\simeq-2\alpha/h$。
相对于外大气，水层使每个圆盘承受
朝向另一盘的净压强 $2\alpha/h$。
乘以整个圆盘面积 $\pi R^2$：
\[
\boxed{F=\frac{2\pi R^2\alpha}{h}}.
\]
使用水表面张力
$\alpha\simeq0.073\,\mathrm{N/m}$，
\[
\boxed{F\simeq
\frac{2\pi(0.05)^2(0.073)}{1.9\times10^{-6}}
\simeq6.0\times10^2\,\mathrm N=0.60\,\mathrm{kN}}.
\]
\excheck 尽管水层质量可以很小，
但压差乘较大的接触面积仍会产生显著吸附力。

\thermocapillaryproblem{176}{82}
两块部分浸入水中的竖直平行玻璃板，
板间距
$d=0.10\,\mathrm{mm}$、
板宽 $\ell=12\,\mathrm{cm}$。
假定完全润湿，且板间水面
不会达到板的上缘，求它们的相互吸引力。
\exsolution
板间的毛细上升高度为
\[
H=\frac{2\alpha}{\rho gd}
\]
（润湿、板间间距 $d$ 很小）。
相对外侧大气与水面，板间水柱内
高出外水面 $z$ 处的压差大小约为
$\rho g(H-z)$。
每块板受内侧压差作用的面积微元
$\ell\,\mathrm dz$，
故两板间吸引力大小为
\[
F=\ell\int_0^H\rho g(H-z)\,\mathrm dz
=\frac12\rho g\ell H^2.
\]
代入 $H$：
\[
\boxed{F=\frac{2\alpha^2\ell}{\rho gd^2}}.
\]
以 $\alpha\simeq0.073\,\mathrm{N/m}$
代入题目数据，
\[
\boxed{F\simeq13\,\mathrm N}.
\]
\excheck 需有足够高的板面容纳毛细上升高度；
本题特别给出这一几何条件。

\thermocapillaryproblem{177}{82}
半径为 $R$ 的肥皂泡
通过长为 $\ell$、内半径为 $r$ 的细管
与大气连通。
肥皂水表面张力为 $\alpha$，
气体动力黏度为 $\eta$。
求肥皂泡排尽气体而消失的时间。
\exsolution
肥皂泡有内外两层膜，
其内部超压为
$\Delta p=4\alpha/R_b$，
其中 $R_b(t)$ 为泡的瞬时半径。
按泊肃叶定律，细管气体的体积流量
\[
Q=\frac{\pi r^4\Delta p}{8\eta\ell}
=\frac{\pi\alpha r^4}{2\eta\ell R_b}.
\]
在近似不可压缩的低超压流量模型中，
泡体积
$V_b=4\pi R_b^3/3$，
且 $-\mathrm dV_b/\mathrm dt=Q$。
因此
\[
-4\pi R_b^2\frac{\mathrm dR_b}{\mathrm dt}
=\frac{\pi\alpha r^4}{2\eta\ell R_b},
\quad
\mathrm dt=-\frac{8\eta\ell}{\alpha r^4}R_b^3\,\mathrm dR_b.
\]
从初始 $R_b=R$ 积分至0：
\[
\boxed{\tau=\frac{8\eta\ell}{\alpha r^4}
\int_0^R R_b^3\,\mathrm dR_b
=\frac{2\eta\ell R^4}{\alpha r^4}}.
\]
\excheck 忽略气体可压缩性和泡膜的破裂机制；
真实泡在到达数学上的零半径前可能失稳破裂。

\thermocapillaryproblem{178}{82}
把竖直毛细管下端接触水面，
水在管内自发上升。
假定完全润湿，水表面张力为 $\alpha$。
求上升过程放出的热量 $Q_{\rm out}$。
\exsolution
设毛细管内半径为 $r$，
平衡上升高度
$H=2\alpha/(\rho gr)$。
毛细上升时毛细力可做的总功为
$W_\alpha=2\pi r\alpha H$。
最终液柱重力势能为
\[
U_g=\int_0^H\rho g(\pi r^2)z\,\mathrm dz
=\frac12\rho g\pi r^2H^2.
\]
由平衡条件
$\rho g\pi r^2H=2\pi r\alpha$，
所以
$U_g=\tfrac12W_\alpha$。
忽略动能耗散后的最终储能，
其余一半功以热释放：
\[
Q_{\rm out}=W_\alpha-U_g
=\frac12\rho g\pi r^2H^2.
\]
代入 $H$：
\[
\boxed{Q_{\rm out}=\frac{2\pi\alpha^2}{\rho g}}.
\]
\excheck 在该理想模型中结果与毛细管半径无关；
这反映更细的管上升更高，
两种半径依赖恰好抵消。

\thermocapillaryproblem{179}{83}
求液体表面层的自由能：
(a) 直径 $d=1.4\,\mathrm{mm}$ 的水银滴；
(b) 直径 $d=6.0\,\mathrm{mm}$ 的肥皂泡，
其肥皂液表面张力为
$\alpha=45\,\mathrm{mN/m}$。
\exsolution
自由表面面积为 $S$，
表面层自由能在给定温度下为
$F=\alpha S$。
(a) 一颗球形水银滴只有一个外表面，
$S=4\pi R^2=\pi d^2$，所以
\[
\boxed{F_{\rm Hg}=\pi\alpha_{\rm Hg}d^2}.
\]
取水银
$\alpha_{\rm Hg}\simeq0.48\,\mathrm{N/m}$，
得
$\boxed{F_{\rm Hg}\simeq3.0\,\mathrm{\mu J}}$。
(b) 肥皂泡有两个气液表面，
面积总和 $2\pi d^2$，故
\[
\boxed{F_{\rm soap}=2\pi\alpha d^2}.
\]
取 $d=0.006\,\mathrm m$ 与
$\alpha=0.045\,\mathrm{N/m}$，
得
$\boxed{F_{\rm soap}\simeq10\,\mathrm{\mu J}}$。
\excheck 肥皂泡的两层膜面须乘2，
但水银滴只乘一个表面面积。
